%%%%%%%%%%%%%%%%%%%%%%%%%%%%%%%%%%%%%%%%%%%%%%%%%%%%%%%%%%%%%%%%%%%%%%%%%%%%%%%
%% Mono-Width Parametric Bitvectors (PBV_1) --- based on the OOPSLA'25 paper.
%% Shared macros and generated plot data live in
%% chapters/parametric-bv/pbv-preamble.tex, which thesis.tex inputs once before
%% this chapter and the multi-width chapter.
%%%%%%%%%%%%%%%%%%%%%%%%%%%%%%%%%%%%%%%%%%%%%%%%%%%%%%%%%%%%%%%%%%%%%%%%%%%%%%%
\chapter{Mono-Width Parametric Bitvectors}
\label{ch:mono-width}
\epigraph{The limits of my language mean the limits of my world}{Ludwig Wittgenstein, \emph{Tractatus Logico-Philosophicus}}

\paragraph*{Attribution.} This chapter is based on \emph{Certified Decision
Procedures for Width-Independent Bitvector Predicates}~\citep{bhat2025certified},
published at OOPSLA 2025, joint work with Léo Stefanesco, Chris Hughes, and
Tobias Grosser.
\emph{My contributions:} the initial algorithm, of handling the mono-width case and
equality predicates only, was developed by Chris Hughes. I noticed its
connection to Presburger arithmetic, reinterpreted the classical decision
procedure as an inductive ($k$-induction) safety proof, implemented the verified
procedure in Lean together with its tactic-preprocessing and reflection
frontends, and generalized it to decide all predicates of the theory; I also
implemented the mixed-Boolean-arithmetic solver. Léo Stefanesco contributed an
alternative solver based on automatic structures.

\section{Overview}
\label{sec:swbv:overview}
% Two paragraphs, essentially no citations.
% P1: the question this chapter answers, where it sits in the thesis argument,
%     and the result in one plain sentence.
% P2: the approximations made to the theoretically-justified algorithm, then
%     the empirical payoff.

% underpin automated reasoning about compilers, hardware,
% and cryptography, yet their \emph{widths} are as much a part of the problem as
% their values. Fixed-width solvers must know every width in advance, while a
% great many real-world goals, such as compiler rewrites,
% require us to prove an identity for \emph{all} widths at once.
% This chapter takes the first step toward deciding such \emph{parametric} bitvector goals: it decides
% predicates over a \emph{single} symbolic width, called the \emph{mono-width} fragment,
% $\text{PBV}_1$. It does so by casting them as model-checking problems and mechanizing the
% requisite automata and $k$-induction machinery.

\lettrine{B}{itvectors} are foundational for automated reasoning. A few
interactive theorem provers, such as Lean, have strong support for deciding
\emph{fixed-width} bitvector predicates by means of bitblasting
(\cref{sec:bg:bvdecide}). However, even these ITPs provide little automation
for \emph{width-independent} bitvector predicates: goals that must hold for
\emph{all} bitwidths, such as the rewrite rules a compiler applies to
integers of every size, or the normalization lemmas that fixed-width
bitblasters themselves rely on~\citep{niemetz2024scalable}. Fixed-width
solvers also face scaling issues as widths grow, struggling beyond width 512
in domains such as hardware design and
cryptography~\citep{lopes2018future,hydra}.

To fill this gap, we take up an old observation. Classical algorithms for
fragments of the width-independent theory can be viewed through the lens of
\emph{model checking}: the formula corresponds to an automaton, and its
correctness is a safety property. But this lens was previously unusable in
mechanized proofs, since there were no executable, fast, formally verified
model-checking algorithms usable \emph{from within} an ITP. We therefore
mechanize key algorithms from the model checking literature: $k$-induction,
automata reachability, automata emptiness checking, and automata
minimization, and use them to build scalable, certified decision procedures.
For the controlled mixtures of arithmetic and bitwise operations that occur in
the deobfuscation literature, we further mechanize the fast MBA-Blast
algorithm, which outperforms the general procedures on that fragment.

\section{Notation \& Background}
\label{sec:swbv:background}

\paragraph{Notation} $\N$ are the natural numbers. We assume an infinite set of variables $\V$.
$\B$ is the finite field of Booleans.
%
We define bitstreams as infinite sequences of booleans ($\bitstream \eqdef \N
\to \B$).

\paragraph{Bitvectors}
A bitvector~$x$ of width~$w \in \N$ is a sequence of $w$ booleans, and the
corresponding sequence of bits is written $\langle x_{w-1}, \ldots, x_0
\rangle$. For example, $x = \langle 1, 0, 1 \rangle$ is a bitvector of length 3.
%
We write $\BV{w}$ for the set of bitvectors of width~$w$.
%
Every bitvector~$x$ can be interpreted as a natural number, by interpreting the
sequence of bits as a binary number $\natofbv{x} \eqdef \sum_{i=0}^{w-1} x_i *
2^i$ (In this case, $\natofbv{x} = 1 \cdot 2^2 + 0 \cdot 2^1 + 1 \cdot 2^0 =
5$). We can also interpret a bitvector as a two's complement integer,
$\intofbv{x} \equiv \sum_{i=0}^{w-2} x_i * 2^i - x_{w-1} * 2^{w-1}$ (In this
case, $\intofbv{x} = 1 \cdot 2^1 + 0 \cdot 2^0 - 1 \cdot 2^2 = -3$). Bitvectors
are both a boolean algebra (a sequence of bits) and a ring (seen as integers
modulo~$2^w$), and the two algebraic structures mix in subtle ways when reasoning
about bitvectors (e.g., $-x = \lnot x + 1$).

\paragraph{Signed and Unsigned Comparison}
Because bitvectors can either be seen as representing signed or unsigned
integers, they are equipped with two relevant orders depending on how they are
interpreted, and which our decision procedures need to handle:
\begin{align*}
  \textit{Unsigned comparison:}&\qquad x <_u y \;\iff\; \natofbv x < \natofbv y \\
  \textit{Signed comparison:}&\qquad x <_s y \;\iff\; \intofbv x < \intofbv y
\end{align*}

\section{Reducing Width-Independent Bitvector Formulas to Model Checking}
\label{sec:swbv:reduction}

We formally define the fragment of bitvector theory we decide,
and then describe the reduction to automata, followed by the model checking algorithm and its mechanization.

\subsection{Syntax and Semantics of formulas}
\label{sec:swbv:formulas}

\begin{figure}[tb]
\centering
{\footnotesize
\begin{tabular}{llll}
$\mathsf{E}$ & $::=$ & $\ofNat~(v : \N)$ & width-independent constant $v$ \\
    & $\mid$ & $\Var~(x : \V)$ & bitvector variable \\
    & $\mid$ & $\mathsf{E}_1 \boxplus \mathsf{E}_2$ & binary operation, $\boxplus \in \{+, -, \bitor, \bitand, \bitxor\}$ \\
    & $\mid$ & $- \mathsf{E}$ & arithmetic negation \\
    & $\mid$ & $\bitnot \mathsf{E}$ & bitwise complement \\
    & $\mid$ & $\mathsf{E} \ll k$ & left shift by a constant $k \in \N$ \\[4pt]
$\mathsf{F}$ & $::=$ & $\mathsf{E}_1 \otimes \mathsf{E}_2$ & comparison, $\otimes \in \{<_u, \leq_u, \leq_s, <_s, =, \neq \}$ \\
    & $\mid$ & $\width~(r : \WidthRel)~(k \in \N)$ & assertion $w \mathrel{r} k$ on the bitwidth $w$ \\
    & $\mid$ & $\mathsf{F}_1 \land \mathsf{F}_2$ & conjunction \\
    & $\mid$ & $\mathsf{F}_1 \lor \mathsf{F}_2$ & disjunction \\
    & $\mid$ & $\mathsf{F}_1 \Rightarrow \mathsf{F}_2$ & implication \\
    & $\mid$ & $\neg \mathsf{F}$ & negation \\[4pt]
$\WidthRel$ & $::=$ & ${<} \mid {\leq} \mid {>} \mid {\geq} \mid {=} \mid {\neq}$ & relations on the bitwidth \\
\end{tabular}
}
  \caption{Syntax of the width-independent bitvector theory: expressions
  $\mathsf{E}$ over a single symbolic bitwidth $w$, formulas $\mathsf{F}$, and
  relations $\WidthRel$ for width assertions.}
  \label{fig:swbv:syntax-formula}
\end{figure}

\paragraph{Syntax}
The syntax of bitvectors and formulas is given in
Figure~\ref{fig:swbv:syntax-formula}.
%
All the bitvectors have the same symbolic bitwidth which we denote by~$w$ in
this chapter.
%
The base expressions are constants ($\ofNat~v$) and variables ($\Var~i$).
The constant expression $(\ofNat~v)$ denotes a \emph{width-independent}
constant: For example, $(\mathsf{ofNat}~3)$ denotes the empty bitvector $\langle
\rangle$ at width 0,  $\langle1\rangle$ at width 1, $\langle 11 \rangle$ at
width 2, $\langle 011 \rangle$ at width 3, and so on. These atoms can be
combined using the usual bitwise operations (bitwise or~$\bitor$, and~$\bitand$,
xor~$\bitxor$, negation~$\bitnot$ as well as left shift $\ll$ by a constant) and
arithmetic addition, subtraction and negation.

The formulas are built from relations for equality of bitvectors, as well as
signed and unsigned inequalities, which can be combined using the usual Boolean
operators.
%
Moreover, we support predicates which are assertions about the widths of the bitvectors.
For instance, $\width~{(\leq)}~4$ states that the width $w$
is less than or equal to~4. This can, for example, be useful to prove results
that only hold for non-trivial bitwidth~($w > 0$).
%
Since the atomic predicates (equality, signed and unsigned comparison) are
closed under negation, predicates can be normalized into negation normal form.
%
We write $\FV(\mathsf F)$ the set of free variables of the formula~$\mathsf F$.


\paragraph{Semantics}
The semantics of expressions and formulas are parameterized by a bitwidth~$w$
and \emph{environment}~$\rho$ mapping the free variables of~$E$ to bitvectors of width~$w$.
%
The semantics of an expression $\mathsf{E}$ is written $\bvsemgen{\mathsf{E}}$
and is defined in the standard way:
%
\begin{align*}
  \bvsemgen{\mathsf{\ofNat~v}} \;&\eqdef\; \left< v_{w-1}, \dots, v_0 \right> \\
  &\quad \text{if~$n$ can be written in binary $v_N...v_0$ with $N \geq w-1$}\\
  \bvsemgen{\mathsf{\Var~x}} \;&\eqdef\; \rho(x)\\
  \bvsemgen{\mathsf{E}_1 \boxplus \mathsf{E}_2} \;&\eqdef\;
    \bvsemgen{\mathsf{E}_1} \boxplus \bvsemgen{\mathsf{E}_2}
\end{align*}

\noindent
Similarly, formulas are interpreted as predicates over the bitwidth and the
environment in the obvious way.
%
We can now formally state the decision problem: given a formula~$\mathsf F$,
determine whether
\[
  \forall w \in \N,\; \forall \rho \in \FV(\mathsf F) \to \BV w,\;\; \bvsemgen{\mathsf F} \text{ holds}.
\]

\subsection{Reduction to Automata}
\label{sec:swbv:reduction-to-automata}

The first order theory of bitvectors described in the previous section can be
decided by reducing it to the emptiness problem of a regular language; this was
first discovered by \citet{Büchi1960} for Presburger arithmetic over (unbounded)
natural numbers.
%
Given a formula $\mathsf{P}$ with $n$-free variables~$\vec{x} \equiv (x_1, \dots, x_n)$,
its solutions of
width~$w$ are the subset $\Sol{w}{\mathsf{P}} = \{ \vec\genbv \in (\bv{w})^n \mid
\denoteinline{\mathsf{P}}^w_{\vec\genbv} \}$, and its set of solutions is the set
%
\[
  \Sol{}{\mathsf{P}} \;=\; \bigcup_{w \in \N}\; \Sol{w}{\mathsf{P}}
  \;\;\subseteq\;\;
  \bigcup_{w \in \N}\; (\bv{w})^n
\]
%
The decision problem amounts to deciding whether $\Sol{}{\mathsf{P}}$ is equal
to the whole set of possible solutions $\bigcup_{w \in \N}\; (\bv{w})^n$.

%
This set of solutions is not a \emph{language} itself, as it not composed of
words over a finite alphabet, but is isomorphic to one: a bitvector $\genbv$ of
width~$w$ can be encoded as a word of length~$w$ over the alphabet $\{0,1\}$
using a big endian convention, and a tuple $\vec\genbv$ of $n$~bitvectors of
width~$w$ can be encoded as a word over the alphabet $\{0, 1\}^n$, where the $k$th letter
corresponds to the $k$th bit of each bitvector in~$\vec{\genbv}$.
%
In summary, the solutions (over all widths) of a formula $\mathsf{P}$ can simply
be seen as a language over the alphabet~$\{0,1\}^n$.
%
This representation also has the advantage of being much simpler than the one
above, as the width of the bitvectors is now implicit.

It now remains to see that the language $\Sol{}{\mathsf{P}}$ is regular, and is
therefore recognizable by a finite automaton.
%
As the logical connectives in Figure~\ref{fig:swbv:syntax-formula} correspond to union
and intersection of the solutions, they preserve the regularity of the language
of solutions, and it suffices to recognize the atomic predicates.
%
This can be further split into two problems, that of recognizing terms and relations.

\paragraph{Terms}
We recognize terms $\mathsf{T}$, by providing a finite
automaton~$\aut_{\mathsf{T}}$ over the alphabet~$\{0,1\}^{n+1}$ where $n$ is
the number of free variables in~$\mathsf{T}$. The first $n$ bits of the
letters correspond to the values of the free variables, and the last bit
corresponds to the value of~$\mathsf{T}$. In other words, $\aut_{\mathsf{T}}$
is a transducer. For example, Figure~\ref{fig:swbv:addition-nfa} implements the XOR
operation and addition:
%
For XOR, there is single state, and the transitions $[x_1, x_2, y]$
express $y = x_1 \bitxor x_2$ for the 4 possible values of the
inputs~$x_1$ and~$x_2$.
%
For addition, the states of the NFA correspond to the carry, and each
transition labeled with $[x_1, x_2, y]$ determines the output bit~$y$ and the
new carry state based on the two input bits $x_1$ and~$x_2$ and the current
carry state.
%
We provide an automaton for each of the constructors of
expressions~$E$ in Figure~\ref{fig:swbv:syntax-formula}.

Of course, this technique is limited by the expressivity of NFAs.
%
For example, multiplication cannot be represented by a finite automaton
because the $i$th output bit cannot be determined from the $i$th input bit and
a finite state.
%
Another, perhaps more surprising, non-example are \emph{right} bitshifts by a
constant~$k$: here the issue is that the $i$th bit of the output depends on the
$i+k$th bit of the input, so the automaton needs a delay of~$k$ bits before
starting to output bits.
%
We believe this is not a fundamental limitation of NFAs and, if we considered
a more complex encoding of bitvectors into strings with a special letter for
delay, we could handle right shifts.

\paragraph{Relations}
We recognize each basic relation $\otimes \in \{<_u, \leq_u, \leq_s, <_s,
=, \neq \}$, by providing an automaton~$\aut_\otimes$ over the
  alphabet~$\{0,1\}^2$ which recognizes the relation~$\otimes$, in the sense
  that $\genbv_1 \otimes \genbv_2$ iff the automaton $\aut_\otimes$ recognizes the encoding
  of the pair~$(\genbv_1, \genbv_2)$.
%
For instance, the automaton for unsigned comparison $\leq_u$ has three states,
corresponding to the following cases:
\begin{enumerate}
\item the state $=$ corresponds to the case where, so far, the bits of the two
  bitvectors are equal;
\item the state $<$ corresponds to the case where, based on the bits that have
  been read, the first operand is smaller;
\item the state $>$ is dual to~$<$.
\end{enumerate}
%
The initial state is~$=$, and the accepting states are $=$ and $<$ since we
wish to recognize~$\leq_u$.
%
The automaton for signed comparison~$\leq_s$ is similar, but has two additional
states $<_\textsf{end}$ and $>_\textsf{end}$ without outgoing transitions, so
that they can only be reached by the last letter, and allows to treat the
sign bit differently from all the others. The accepting states are $=$ and
$<_\textsf{end}$.

\begin{figure}
  \usetikzlibrary {arrows.meta,automata,positioning}
  \begin{tikzpicture}[->,>={Stealth[round]},shorten >=1pt,
                      auto,node distance=4cm,on grid,semithick,
                      inner sep=2pt,bend angle=45]
    \node[initial,state, accepting] (X) {};
    \path [every node/.style={font=\footnotesize}]
          (X) edge [loop above] node {\lettervec000 \lettervec011 \lettervec101 \lettervec110} (X);

    \node[initial,state, accepting, right=of X] (Z) {$0$};
    \node[state, accepting] (S) [right=of Z] {$1$};
    \path [every node/.style={font=\footnotesize}]
          (Z) edge [bend left] node {\lettervec110} (S)
          (S) edge [bend left] node[above] {\lettervec001} (Z)
          (S) edge [loop above] node {\lettervec010 \lettervec100 \lettervec111} (S)
          (Z) edge [loop above] node {\lettervec000 \lettervec011 \lettervec101} (Z);
          \end{tikzpicture}
  \caption{Automata implementing XOR (left) and addition of bitvectors (right)}
  \label{fig:swbv:addition-nfa}
\end{figure}

%
Using standard techniques from automata theory, these automata can be combined
into an automaton recognizing any formula.
%
This technique is an instance of a more general framework called \emph{automatic
structures}~\citep{automatic-structures}.

\section{Automata-Based Decision Procedure}
\label{sec:swbv:automata}

The first model checking method presented here uses standard automata theory to
directly decide if the automaton $\aut_{\mathsf{P}}$ is \emph{universal}, that
is, whether it recognizes all the words over its alphabet.
%
At a high level, we first develop an efficient automata library in Lean, and
use it to define
\begin{itemize}
  \item a function \leaninline{nfaOfFormula} that takes a
    formula~$\mathsf{P}$ as input and produces an automaton which recognizes its
    solutions,
  \item a function \leaninline{isUniversal} which returns a Boolean
    indicating whether an automaton is universal.
\end{itemize}
%
By composing them, we obtain a function
\begin{lean}
  def formulaIsUniversal (F : Formula) : Bool :=
    (nfaOfFormula F).isUniversal
\end{lean}
which satisfies the following correctness theorem:

\begin{lean}
  theorem decision_procedure_is_correct (F : Formula) :
      formulaIsUniversal F = true → ∀ (w : Nat) (env : Nat → BitVec w), F.sat env
\end{lean}

This allows us to do \emph{proofs by reflection}: to prove a Lean
formula~$\phi$, a tactic computes a syntactic formula~$\mathsf{F}$ and an
environment $\rho$ such that $\phi \iff \denote{\mathsf{F}}_\rho$.
%
It then suffices to check that the Lean Boolean
term~\leaninline{formulaIsUniversal F} is equal to true.
%
To be more efficient, we ask Lean to compile
\leaninline{formulaIsUniversal} to native code, and to invoke this native
code during type checking.

\subsection{A Certified and Efficient NFA Library}

As part of this decision procedure, we implement a formally verified NFA library
in Lean which needs to be efficient in order to decide large problems.

\paragraph{Data Representation}

One important feature for efficiency is that the states of the automaton do not
have an identity, they are instead natural numbers between 0 and
\leaninline{stateMax}. Thus, a \emph{computational} NFA over the alphabet
\leaninline{A} is defined as follows in Lean:

\begin{lean}
  structure RawCNFA (A : Type) where
    stateMax : Nat
    initials : Std.HashSet Nat
    finals   : Std.HashSet Nat
    trans    : Std.HashMap (Nat × A) (Std.HashSet Nat)
\end{lean}

To be able to reason about such automata, we define a well-formedness predicate
\leaninline{m.WF} which states that all the integers in the
\leaninline{RawCNFA} are less than~\leaninline{m.stateMax}.
%
A \leaninline{CNFA} is then a (dependent) pair of a \leaninline{RawCNFA} and a
proof that it is well-formed.
%
We found it easier to work with a simply typed raw data-structure with an
external well-formedness predicate than with an intrinsically typed version that
would use dependent pairs \leaninline{{s : Nat // s < stateMax}} as states
instead of simple integers, as adding a new state to an automaton would
increment \leaninline{stateMax} and thus invalidate this type.
%
All high-level operations are exposed in terms of
\leaninline{CNFA}.

\paragraph{Operations over Automata}

There are a handful of widely useful operations over NFAs, namely the product
construction (to define disjunction and conjunction of formulas) and the
powerset construction (to determinize, and thus to negate).
%
In the usual textbook description of these operations, the states of the
resulting automaton are (respectively) pairs of states of the input automatons,
and subsets of the input operation. This works well when the goal is to reason
about these constructions, but it is computationally inefficient, as
unreachable states are constructed.

A more efficient method is to use a \emph{worklist} algorithm, where states to
be processed are put on a queue and processed until it becomes empty, as
described for example by \citet{esparza2023automata}.
%
However, termination of worklist algorithms is non-trivial to establish
formally, as one needs to show that the (finite) set of states of the resulting
automaton that have not yet been explored is decreasing.
%
To amortize this effort, we define a generic worklist algorithm for NFAs which
dramatically reduces the required proof effort for the individual operations.

\paragraph{Generic Construction Algorithm}
The idea is that, although our automata states do not have identities,
\emph{during the construction}, we need to keep track of the mapping between
states in the new automaton and state(s) in the input automata.
%
Consider, for instance, the product construction: suppose we want to compute the
product automaton of \leaninline{m1 : CNFA A} and \leaninline{m2 : CNFA A}.
%
During the construction of the product, given a state~$s_{\mathit{new}}$ of the
product automaton, we need to know what the successor states of
$s_{\mathit{new}}$ should be, since we are building the automaton incrementally
starting from its initial states and following transitions.
%
For this, we need to know that it corresponds to some pair $(s_1, s_2)$ of the
input automata \leaninline{m1} and \leaninline{m2} in order to be able to
determine its successor states and potentially add new states to the worklist:
states that correspond to pairs $(s'_1, s'_2)$ reachable from $s_1$ and
$s_2$ using the same letter.
%
Generalizing, we posit a type of states~\leaninline{S} together with operations:
%
\begin{lean}
  step  : S → Array (A x S)
  inits : Array S
  final : S → Bool
\end{lean}

We think of the elements of~\leaninline{S} as the states of the automaton under
construction: for example pairs of states (integers) for the product or finite
sets of states of the input automaton for the powerset construction.
%
The worklist algorithm keeps track of the \emph{worklist}, the set of elements
of \leaninline{S} that need to be processed and a \emph{mapping} from
\leaninline{S} to integers corresponding to the states of the new automaton.
%
The initial state of the queue is given by \leaninline{inits}, and the outgoing
transitions of a state~$s$ are given by~\leaninline{step s}, and the new states
are suitably added to the worklist. Finally, the \leaninline{final} function
determines which states are final.
%
Termination of this algorithm is proved once and for all by establishing that
the set of states in \leaninline{S} that are not in the mapping or are in the
worklist is strictly decreasing.
%
With this generic algorithm at hand, defining these constructions essentially
boils down to writing down their textbook definitions.

\paragraph{Specification by Bisimulation}

We choose to never directly reason about the language accepted by a
\leaninline{CNFA}, instead we use Mathlib's \citep{mathlib} \leaninline{NFA} type, which
corresponds to the textbook definition of an automaton as a tuple $\left<Q, A,
I, F, \delta\right>$ of a set~$Q$ of states, the subsets $I$ and~$F$ of initial
and final states, and a transition relation~$\delta$.
%
We then completely specify a \leaninline{CNFA} using an abstract NFA to
which it is \emph{bisimilar}~\cite{kozen-automata-and-computability}.
%
This is an extension of the notion of bisimulation and of bisimulation from
labeled transition systems (LTS) to NFAs: two NFAs $\aut_1$ and~$\aut_2$ over
the same alphabet are \emph{bisimilar} if there exists a relation~$R$ between
their sets of states such that:
\begin{enumerate}
  \item forgetting initial and final states, $R$ is a simulation between the
    LTSs $(Q_1, \delta_1)$ and $(Q_2, \delta_2)$, that is, if $R(q_1, q_2)$ and
    $q_1' \in \delta_1(q_1, a)$, then there exists $q_2' \in \delta_2(q_2, a)$
    such that $R(q_1', q_2')$;
  \item for any initial state~$q_1 \in I_1$ of $\aut_1$, there exists an initial
    state $q_2$ of $\aut_2$ such that~$R(q_1, q_2)$, and vice-versa;
  \item given $R(q_1, q_2)$, $q_1$ is final in~$\aut_1$ iff $q_2$ is final in~$\aut_2$.
\end{enumerate}

The choice of NFA as specification does not matter, since two bisimilar
automata necessarily recognize the same language:
%
The advantage of this strategy is that, once the bisimulation has been
established, we can reason about the properties of mathematical NFAs, a
substantially easier task.

One example of this style of specification is the generic worklist algorithm:
with only very weak assumptions on~\leaninline{S} and its associated
operations (to ensure termination), we establish that the
\leaninline{CNFA} it produces is in bisimulation with the NFA
%
\[
  \aut_W \;\eqdef\;
  \left<\text{\leaninline{S}}, \text{\leaninline{A}},
  \text{\leaninline{inits}}, \{q \mid  \text{\leaninline{final}}~q =
  \text{\leaninline{true}}
  \} , \text{\leaninline{step}} \right>
\]
%
It then suffices to provide a bisimulation between the NFA above and, say, the
product automaton, which is straightforward.
%
For instance, when constructing the product automaton, we assume that we
have~\leaninline{m1} and \leaninline{m2} which are (respectively) bisimilar to
two abstract NFAs $\aut_1$ and $\aut_2$, with relations~$R_1$ and $R_2$ between
their states.
%
When running the worklist algorithm, we have chosen $S = \N \times \N$
(corresponding to states of the two input automata). It is then straightforward to
prove that the relation $R$ defined by
\[
  R((s_1, s_2), (q_1, q_2)) \;\iff\; R_1(s_1, q_1) \;\land\; R_2(s_2, q_2)
\]
is a bisimulation between $\aut_W$ and $\aut_1 \times \aut_2$, by proving that
the parameters \leaninline{inits}, \leaninline{step} and \leaninline{final}
passed to the worklist algorithm correspond to the standard product construction.

Using bisimilarity, we define the language (formalized as a set of lists of
letters) recognized by a CNFA as follows:
\begin{lean}
  def CNFA.recognizes (m : CNFA A) (L : Language A) :=
    ∃ (M : NFA A), m.Sim M ∧ M.accepts = L
\end{lean}
and prove that this relation is functional (in the sense that, if an automaton
recognizes $L_1$ and $L_2$, then $L_1 = L_2$).

\paragraph{Minimization} To minimize an automaton~$\aut$, we use Brzozowski's algorithm
as it is simple to implement and verify, it consists of reversing (reversing all
the transitions, and swapping initial and final states) and
determinizing the NFA twice, yielding the minimal deterministic automaton
recognizing the same language as~$\aut$.

One difficulty is that simulation and reversal are not compatible, in the sense
that it is not true that, if $\aut_1$ is bisimilar to $\aut_2$, then $\aut_1^R$
is bisimilar to $\aut_2^R$. Indeed, the notion of (forward) simulation is
fundamentally oriented, and does not say anything about backwards simulation.
%
There is a trick though: Since we only care about the language recognized
by~$\aut_2$, we can replace~$\aut_2$ with another automaton which is easier to
work with, as long as it is bisimilar to~$\aut_1$, namely $\aut_1$ itself! The
proof that the language of~$\aut^R$ is the reverse of the language of~$\aut$
becomes easy, and thus the proof that Brzozowski's algorithm is correct.


\subsection{Mechanization of the Reduction to Automata}

With this NFA library in hand, it is relatively straightforward to implement the
encoding of bitvector formulas into automata:

\paragraph{Correspondence Between Solutions and Language}

We formalize the isomorphism between the set of solutions of the
formula, which is a set of vectors of $n$ bitvectors of arbitrary bitwidths,
and the language over the alphabet \leaninline{BitVec n}, by writing an encoding
and decoding function, and proving they are inverses of each other.
%
The former is defined as a dependent pair of a number~$w$ and a vector of~$n$
bitvectors of width~$w$:
%
\begin{lean}
 structure BitVecs (n : Nat) where
  w : Nat
  bvs : Vector (BitVec w) n
\end{lean}
%
We then give a specification to each construction on NFAs over the alphabet~$\BV
n$ in terms of the set of \leaninline{BitVecs n} that it accepts.

\paragraph{Composition of Relations and Operations}

To simplify the encoding, we use an over-approximation of the set~$\FV(\mathsf
F)$ of free variables: we only keep track of $n = \max \FV(\mathsf F)$.
%
Each formula~$\textsf{F}$ is then interpreted as an NFA over the alphabet $\BV
n$. As such, when defining the automaton for, say, $\mathsf F_1 \land \mathsf
F_2$, the product construction cannot be used directly, as the two
automata~$\aut_1$ and~$\aut_2$ have (in general) different alphabets, when the
two formulas have different free variables.
%

The solution is to \emph{lift} both automata to a common alphabet: given an
automaton $\aut$ over the alphabet~$A^n$, and a function $f : \{0, \dots, n-1\}
\to \{0, \dots, m-1\}$, the lifting of~$\aut$ along~$f$ is an automaton over the
alphabet $A^m$ with the same states, and whose transitions are defined as $q
\xrightarrow{\vec{a}} q'$ whenever $q \xrightarrow{f^*(\vec{a})} q'$ in~$\aut$.
%
The $i$th component of the vector $f^*(\vec{a})$ is defined to be equal
to~$a_{f(i)}$.

For expression, the situation is similar, if more complex: an expression
with~$n$ free variables is interpreted as an automaton over the
alphabet~$\BV{(n+1)}$, where the last component corresponds to one bit of the
\emph{result} of the expression.
%
Therefore, to interpret $\mathsf E_1 + \mathsf E_2$, we need to build an
automaton over~$\BV{(k+3)}$, with $k = \max(n_1, n_2)$.
%
The bit~$k$ corresponds to the result of~$\mathsf E_1$, the bit~$k+1$ to that of
$\mathsf E_2$ and the last bit, at position $k+2$, is the result of the addition
of the bits~$k$ and~$k+1$.
%
We then remove the bits~$k$ and~$k+1$ using \emph{projection}, an operation dual
to \emph{lifting}.

With these tools at hand, it is straightforward to interpret any formula as an
NFA. We choose to apply a minimization procedure at each step of the process
in order to keep their size manageable.

\paragraph{Automata for the Base Cases}

The standard constructions on NFAs handle the Boolean structure of formulas.
%
To handle the operations on terms and the base logical predicates over
bitvectors, we construct explicit NFAs which recognize them, this (somewhat
tedious) work is done using the following procedure:
\begin{enumerate}
  \item write a program constructing \leaninline{m : CNFA}, it is usually
    straightforward;
  \item describe the same automaton as an abstract \leaninline{NFA}~$\aut$ with
    a concrete set $Q$ of states;
  \item prove there is a simulation between \leaninline{m} and $\aut$, by
    defining a bijection from~$Q$ to $\{0, ..., N-1\}$ where $N$ is the
    (concrete) size of~\leaninline{m};
  \item prove $\aut$ recognizes the right language, by providing, for each~$q
    \in Q$ the language it recognizes and checking that each transition respects
    this assignment.
\end{enumerate}
%
Since these constructions are bounded with a concrete size (except for the width
predicates), \emph{ad hoc} proof automation is usable and the proof effort is
moderate.

\section{Model Checking by \texorpdfstring{$k$}{k}-Induction}
\label{sec:swbv:kinduction}

The $k$-induction based decision procedure uses a similar construction to the
previously described automata encodings, but instead, chooses to view the
automata as \emph{transducers}, whose outputs are infinite bitstreams.

\begin{figure}[t]
\centering
\small
\setlength{\tabcolsep}{3pt}
\renewcommand{\arraystretch}{1.2}
\begin{tabular}{@{}p{12.5cm}@{}}
\toprule
\textbf{Bitstream semantics} (relative to an environment~$\rho$) \\
\midrule
\textit{Terms}: $\denote{-} : \mathsf{Term} \times \mathsf{Env} \to (\N \to \B)$, \quad where $\mathsf{Env} \equiv \V \to (\N \to \B)$ \\
\dimtext{every term stream is $0$ at position $i = 0$; the rules below give position $i > 0$} \\[2pt]
\rowcolor{black!10!white}
$\denote{\Var(x)}_\rho(i) \defn \rho(x)(i)$ \\
$\denote{\ofNat(v)}_\rho(i) \defn \lfloor v / 2^i \rfloor \bmod 2$ \\
\rowcolor{black!10!white}
$\denote{s \bitor t}_\rho(i) \defn \denote{s}_\rho(i) \bitor \denote{t}_\rho(i)$ \\
$\denote{\lnot s}_\rho(i) \defn 1 - \denote{s}_\rho(i)$ \\[4pt]
\midrule
\textit{Formulas}: $\denote{P}_\rho(i) = 1$ iff $P$ holds at bitwidth $i$ for environment $\rho$ \\[2pt]
\rowcolor{black!10!white}
$\denote{s = t}_\rho(i) \defn \forall j \leq i,\; \denote{s}_\rho(j) = \denote{t}_\rho(j)$ \\
$\denote{P \land Q}_\rho(i) \defn \denote{P}_\rho(i) \land \denote{Q}_\rho(i)$ \\
\bottomrule
\end{tabular}
\caption{Bitstream semantics of the bitvector theory, relative to an
environment~$\rho$ that maps each variable to a bitstream. Terms are
interpreted as bitstreams, extending the bitwise and arithmetic operations
pointwise; a formula is interpreted as the non-increasing bitstream whose
$i$th bit records whether the formula holds at bitwidth~$i$.}
\label{fig:swbv:stx-sem-bv-theory}
\end{figure}

\subsection{Reduction to Transducers}

The problem of deciding a bitvector formula with a universally quantified
bitwidth is reduced to checking that the same formula interpreted in the domain
of infinite bitstreams always outputs~$1$.
%
The interpretation of the terms~$\mathsf E$ and
the formulas~$\mathsf F$ (defined in Figure~\ref{fig:swbv:syntax-formula}) in terms
of bitstreams is presented in Figure~\ref{fig:swbv:stx-sem-bv-theory}.
%
The interpretation of terms is the straightforward extension of bitwise and
arithmetic operations to bitstreams. Formulas are interpreted as bitstreams
whose $k$th bit indicates whether the predicates holds for bits at
positions~$\leq k$.
%
As such, the interpretation of a formula is always a non-increasing sequence.

Thus, deciding whether a formula~$\mathsf F$ universally holds boils down to
deciding whether its bitstream interpretation satisfies $\forall w~\rho,
\denote{\mathsf F}_\rho(w) = 1$.
%
That is to say, whether the formula $\mathsf F$ holds for all bitwidths and
environments. The difference from the automata-based approach is that rather
than labeling states as accepting or rejecting, it instead builds
\emph{transducers}, which produce an output bit for each state.

The transducers are abstractly a tuple $(S, s_0 : S, \iota : \N, \pi : S \times
\bv{\iota} \to \B, \delta : S \times \bv{\iota} \to S)$, where $\iota$ is the
arity of the transducer (e.g., constants have arity 0, negation has arity 1,
addition has arity 2), $\pi$ is the output function, $\delta$ is the state
transition function. However, since the set $S$ of states is \emph{finite}, we
can always replace it with $\bv{\sigma}$, where $\sigma$ is large enough to
encode all states of the transducer.
%
The encoding of transducers is now a tuple $(\sigma : \N, s_0 :
\bv{\sigma}, \iota : \N, \pi : \bv{\sigma+\iota} \to \bv{1}, \delta :
\bv{\sigma+\iota} \to \bv{\sigma})$.
%
Next, since the \emph{function spaces} of $\pi$ and $\delta$ are themselves
finite, we can encode $\pi$ and $\delta$ as boolean circuits, which will be
exploited by the $k$-induction algorithm to build a SAT formula.
%
The final definition is the following:
%
\begin{lean}
structure FSM (ι  : Type) where
  σ  : Type  -- state space size
  s0 : ι → Bool -- initial state
  π  : CircuitFamily (σ ⊕ ι) Unit -- output function: BitVec[σ+ι] → BitVec[1]
  δ  : CircuitFamily (σ ⊕ ι) σ -- transition function: BitVec[σ+ι] → BitVec[σ]
\end{lean}


Given the finite state transducer that represents a formula $F$,
we use model checking on the finite state transducer $\F_F$ to establish that the
output function $\pi$ on all states reachable from $s_0$ for all sequences of inputs always produces $\true$.
This in turn establishes that $\denote{F}_\rho(w) = \true$ for all widths,
which proves that $\forall w, F(w)$ is true.
This hinges on the fact that $\forall w, F(w)$ is a safety property for the FSM
$\F_F$, and thus, any counter-example is a finite trace, which can be found by
symbolically exploring the state space of $\F_F$. The efficiency of this
algorithm depends on the encoding of $\F_F$ and the algorithm used to establish
the safety invariant that the output function $\pi$ always emits $\true$.

\subsection{Establishing Safety via \texorpdfstring{$k$}{k}-Induction}

\begin{figure*}[htb]
  \begin{tabular}{p{\textwidth -15pt}}
    \textcolor{gray}{Reachable by path $p$:} \\
    {$
      (s : S) \xtosquig{p} (u : S) \equiv
    \begin{cases}
      s = u & \text{if } p = \langle \rangle \\
      \delta(s, p_0) = t \land t \xtosquig{q} u & \text{if } p = \langle p_0;  q\rangle
    \end{cases}$} \\
    \rowcolor{gray!10}
    \textcolor{gray}{Reachable by path $p$, all intermediate states output \true:} \\
    \rowcolor{gray!10}
    {$
      (s : S) \xtosquig{p}_{\true} (u : S) \equiv
      \begin{cases}
        s = u &  \text{if } p = \langle \rangle \\
        \pi(s, p_0) = \true \land \delta(s, p_0) = t \land t \xtosquig{q}_\true u & \text{if } p = \langle p_0;  q\rangle
      \end{cases}$} \\
    \textcolor{gray}{States reachable in \( k \) steps from \( s_0 \) are safe:} \\
    $\InitPrecond_{k} \equiv \forall (t: S)~(i : \B)~(p : \BV{k}), s_0 \xtosquig{p} t \implies \pi(t, i) = \true$ \\
    \rowcolor{gray!10}
    \textcolor{gray}{Safe reachability in \( k \) steps can be extended to \( k+1 \) steps:} \\
    \rowcolor{gray!10}
    $\Ind_{k} \equiv \bigvee_{j=1}^k \forall (s~t : S)~(i : \B)~(p : \BV{k}), s \xtosquig{p}_{\true} t \implies \pi(t, i) = \true$ \\
    \textcolor{gray}{Safety = Preconditions + Inductive Invariant:} \\
    $\Safe_{k} \equiv \InitPrecond_k \land \Ind_k$
  \end{tabular}
\caption{
  We establish the safety property \( \Safe_k \) to show that an FSM always outputs \true{}. \( \InitPrecond_k \) establishes that all states reachable in \( k \) steps from \( s_0 \) are safe, and \( \Ind_k \) allows extending the safety frontier from \( k \) to \( k+1 \).
  Together, these prove the FSM is safe. The procedure either finds a counter-example via \( \InitPrecond_k \) for large \( k \), or establishes \( \Ind_k \), as the set of states satisfying \( \Ind_k \) is strictly increasing.
  See that the final safety property, $\Safe_k$, consists of universally quantified predicates,
  and can thus be proven by showing UNSAT of the negated (existentially quantified) problem.
}
\label{fig:swbv:decision-procedure-pesudocode}
\end{figure*}

We define the safety property $\Safe_k$ as the conjunction of two properties (\cref{fig:swbv:decision-procedure-pesudocode}):
We notate the property of a state $s$ being reachable from $s_0$ by a sequence of inputs $p : \BV{k}$ (for some natural number $k$) as $s_0 \xtosquig{p} s$,
and the fact that $s$ transitions to $t$ with input bit $i$ producing output bit $b$ as $s \xrightarrow{i}_b t$.
$\InitPrecond_k$ establishes that all states reachable in $k$ steps from $s_0$ are safe
and $\Ind_k$ allows us to extend the safety frontier from $k$ steps to $k+1$ steps.
Algorithmically, we either find a counter-example via $\InitPrecond_k$ for a
large enough $k$, or we establish $\InitPrecond_k \land \Ind_k$, which establishes that the FSM is safe for all states.

\paragraph{Soundness}
Assume that we establish $\InitPrecond_k \land \Ind_k$.
So, all states reachable in $k$ steps from $s_0$ are safe.
Also, any state reachable in $k$ steps from $s_0$ can be extended to $k+1$ steps (by the induction hypothesis).
Repeatedly applying this, we see that all states reachable from $s_0$ are safe, since the set of states is finite.
This establishes that the FSM is safe, assuming the algorithm terminates.

\paragraph{Completeness}
The set of states satisfying $\Ind_k$ is strictly increasing, since the set of states satisfying $\Ind_{k+1}$ is at least as large as the set of states satisfying $\Ind_{k}$,
as $\Ind_{k} \implies \Ind_{k+1}$.
However, naive $k$-induction performs induction on the length of paths in the state transition graph of the finite state machine, which is unbounded.
To ensure completeness, we add \emph{simple path constraints},
which express that we check $k$-induction only on simple paths.
This is standard, and prevents the algorithm from getting stuck due to the presence of an unreachable
cycle of good states, that have an out transition into a bad state.

Note that having correctness and termination proves that $k$-induction is \emph{complete} for the bitvector fragment we handle, which guarantees to the user that our automation is robust.

\subsection{Mechanization in Lean}

We mechanize the $k$-decision procedure in Lean,
which uses Lean's inbuilt support for SAT solving to establish the safety property
by producing a SAT formula that is then proven to be unsatisfiable.
Our proof of correctness of the algorithm is to show that $\InitPrecond_k \land \Ind_k$
implies that the FSM produces $\true$ for all inputs by the $k$-induction argument.
The safety of the FSM implies that the corresponding formula is true for all bitwidths.
The proof structure is straightforward, establishes that the circuit we build checks two properties:
(1) Safety of states reachable in $\leq k$ steps from the initial state
(2) Safety on a simple path upto length $k$ implies safety on the $k+1$th step.
These two properties together imply safety for all reachable states.

An interesting point is to compare this to the automata decision procedure ---
While both of these fundamentally rely on exploiting a finite description of the
bitwise behavior of the circuits, the algorithm they employ is quite different.
The automata decision procedure uses the well-known theory of finite automata to
decide formulas, while the circuit based approach instead chooses to take a SAT
inspired view of the same.

\section{Tactic Implementation}
\label{sec:swbv:tactic}

\begin{figure}[htb]
\includegraphics[width=\textwidth]{chapters/single-width-bv/images/k-induction.pdf}
\caption{The structure of the $k$ induction proof by reflection.
See that we go from the predicate, to the bitstream semantics, to the FSM transducer semantics,
which is then proven by the $k$ induction algorithm.}
\label{fig:swbv:k-induction-flow}
\end{figure}

We implement and verify our algorithms in the Lean proof assistant. We
illustrate the overall flow of the algorithm on a Lean-level, universally
quantified statement relating addition and bitwise-or (\leaninline{|||})
(\cref{fig:swbv:k-induction-flow}).
%
The Lean statement is first converted to a statement in terms of negation normal form.
%
Then, it is rewritten into a reflected formula~$\mathsf F$ whose denotation is
the bitvector statement, using \leaninline{Predicate.denote}.
%
\begin{lean}
def Predicate.denote (p : Predicate) (w : Nat) (vars : List (BitVec w)) : Prop :=
  match p with
  | .eq t₁ t₂ => t₁.denote w vars = t₂.denote w vars
  | ... (other cases elided)
\end{lean}
The reflected predicate language is then transformed into the infinite stream semantics (\cref{fig:swbv:stx-sem-bv-theory}),
using \leaninline{Predicate.eval}.

\begin{lean}
def Predicate.eval (p : Predicate) (vars : List BitStream) : BitStream :=
  match p with
  | eq t₁ t₂ => Predicate.evalEq (t₁.eval vars) (t₂.eval vars)
  | ...
\end{lean}

We prove a theorem which shows that the stream semantics agrees with the bitvector semantics.
\begin{lean}
theorem Predicate.eval_eq_denote (w : Nat) (p : Predicate) (vars : List (BitVec w)) :
    p.denote w vars ↔ -- denotation into bitvectors.
    (p.eval vars w = false) -- bitstream semantics.
\end{lean}

Next, the stream semantics of the predicate is made concrete by realizing the bitstream
as the output of a finite state machine. First, we show that every finite state machine denotes an infinite bitstream:

\begin{lean}
/-- `eval p` morally gives the function
  `BitStream → ... → BitStream` represented by FSM `p` -/
def FiniteStateMachine.eval (x : arity → BitStream) : BitStream
\end{lean}

Next, we define that a FSM implements a predicate $p$ iff the stream of outputs
from the FSM agrees with the denotation of the predicate $p$:

\begin{lean}
/-- An `FSMPredicateSolution `t` is an FSM with a witness
    that the FSM evaluates to the same value as `t` does -/
structure FSMPredicateSolution (p : Predicate) extends FSM (Fin p.arity) where
    good : p.eval = toFSM.eval
\end{lean}

We then show that every predicate $p$ has a corresponding FSM that implements it:

\begin{lean}
def predicateToFSM (p : Predicate) → FSMPredicateSolution p
  | .eq t₁ t₂ =>
    let t₁ := termEvalEqFSM t₁ -- Convert term syntax to FSM
    let t₂ := termEvalEqFSM t₂ -- Convert term syntax to FSM
    {
     -- Build equality FSM from term FSMs
     toFSM := fsmEq t₁.toFSM t₂.toFSM
     good := eval_fsmEq_eq_evalFin_Predicate_eq
    }
  | ... (Other predicates, elided)
\end{lean}

With this machinery, we show that each predicate has an FSM associated to it.
This FSM is built using constructs from our library, to build FSMs for terms,
which are combined into FSMs for predicates.
%
As a concrete example, consider the equality predicate:
%
we build an FSM, which first outputs a $\true$, since all bitvectors are equal at width 0,
and then, for each width, checks if the bits are equal.
However, this does not suffice, since given two bitvectors $x \equiv \langle 100\dots\rangle$ and $y \equiv \langle 000\dots \rangle$.
such an FSM would output $\langle \true, \false, \true, \true, \dots \rangle$.
However, we need the FSM to output $\langle \true, \false, \colorbox{blue!5!white}{\false, \false, \dots} \rangle$ after the first bit,
since the bitvectors being disequal at index $1$ implies they are disequal at all widths.
Thus, we \emph{scan} the output of the FSM with boolean and to accumulate the first
difference. In the mechanization, this is:

\begin{lean}
def fsmEq (a : FSM (Fin k)) (b : FSM (Fin l)) : FSM (Fin (k ⊔ l)) :=
  composeUnary FSM.scanAnd (  -- scan with && to accumulate 1st difference.
    (composeUnary (FSM.ls true)  -- At width 0, all BVs are equal.
     composeBinary FSM.eq a b  -- For other widths, check if bit is equal
    ))
\end{lean}

Our verification library has support for composing unary and binary functions over FSMs,
which we use as building blocks to build the composite FSMs for terms and predicates.
The core building block, which allows a composition of $n$-ary FSMs is the \texttt{compose} function.
Given an FSM $p$ of arity $n$, a family of $n$ FSMs $q_i$ of arities $m_i$,
and given yet another arity $m$ such that $m_i ≤ m$ for all $i$,
we can compose $p$ with the $qᵢ$ yielding a single FSM of arity $m$,
such that each FSM $qᵢ$ computes the $i$th bit that is fed to the FSM $p$.

\begin{lean}
def compose (n : Type) (p : FSM n) -- FSM p of arity n
  (ms : n → Type) (q : ∀ (i : n), FSM (ms i)) -- family of FSMs q_i
  (m : Type) (f : ∀ (i : n), ms i → m) : -- m_i ≤ m
  FSM m -- output FSM
\end{lean}

This building block is used to create the functions \leaninline{composeUnary} and \leaninline{composeBinary},
which take a unary or binary function and apply it to the FSMs.
As seen above, we use \leaninline{composeUnary} and \leaninline{composeBinary} to build the FSM for the equality predicate,
and more generally, to build the FSM for the predicate $P$ as a composition of smaller FSMs.
The library has common FSM constructors for (a) lifting bitwise operations on booleans,
(b) concatenating bits to an existing FSM, (c) scanning the output of an FSM with a binary function,
and (d) common arithmetic circuits built from these primitives.

Finally, we show that the FSM for the equality predicate is correct,
by proving that the denotation of the FSM equals the bitstream denotation of the predicate:

\begin{lean}
/-- Evaluation FSM.eq is the same as evaluating Predicate.eq.evalFin. -/
theorem eval_fsmEq_eq_evalFin_Predicate_eq (t₁ t₂ : Term) :
   (fsmEq (termEvalEqFSM t₁).toFSM (termEvalEqFSM t₂).toFSM).eval =
   (Predicate.eq t₁ t₂).evalFin
\end{lean}

See that this establishes a proof of correctness,
which connects the semantics of the finite state machine as a transducer
to the semantics of the term/predicate's denotation as a transducer.

\section{Preprocessing and Normalization of Inputs}
\label{sec:swbv:preprocessing}

The tactic front-end performs heavy preprocessing before performing reflection.
This preprocessing is necessary to ensure that the input formula is in a normal form that our small reflective fragment can consume.
We also use the preprocessing to extend the set of decidable predicates that can be handled by the decision procedure,
to a more expressive, but logically equivalent fragment.
We use Lean's simplifier, \leaninline{simp} to drive our preprocessing.
We register lemmas that simplify the input formula into a normal form.

\paragraph{Normalize Literals.}
The constant bitvector $\langle10\rangle$ can be represented in many ways in Lean:
(a) by the numeral \leaninlineM{3}, which is implicitly cast into a bitvector,
(b) by the integer \leaninlineM{-1 : Int}, which is cast into a bitvector,
(c) by the bitvector \leaninlineM{BitVec.ofNat (w := 2) 3}, which is the explicit bitvector representation.
and (d) by the bitvector expression \leaninlineM{BitVec.ofInt (w := 2) (-1)}, which is the explicit integer representation.
We coerce all of these into \leaninlineM{BitVec.ofNat} as a normal form,
since it has the most complete support for reasoning and constant folding currently in the proof assistant.

\paragraph{Normalize Arithmetic.}
To help with computing normal forms, we perform the usual reassociations and commutations.
First, we move constants to the left of expressions for better constant
folding. This permits the following set of rewrites $1 + (x + 2) \to 1 + (2 + x) \to (1 + 2) + x \to 3 + x$.
Next, we unfold multiplication into a shift-and-add form, so that $x \times 2^l \cdot k$ is
rewritten as $(x \ll l) * k$, and $x \times (2k + 1)$ is rewritten to $x * 2k + x$,
and we rewrite $x \ll l \ll m \to x \ll (l + m)$. These rewrites, taken together,
break a multiplication up into repeated shifts-and-adds. For example,
the rewrites normalize $x * 13 \to x*12 + x \to (x \ll 2) * 3 + x \to (x \ll 2) * 2 + (x \ll 2) + x \to ((x \ll 2) \ll 1) + ((x \ll 2) + x) \to x \ll 3 + x \ll 2 + x$
This allows the core decision procedure to only reason about shifts and adds.

\paragraph{Normalize Predicates.}
Given a goal state of the form $H_1, H_2 \dots \vdash P \to Q$,
we transform the goal into $H_1 \to H_2 \to \dots \to P \to Q$ by using Lean's \leaninline{revert} tactic to revert all hypotheses.
Next, we convert the goal into $\vdash \lnot H_1 \lor \lnot H_2 \lor \dots \lor \lnot P \lor Q$.
This is now in a form that can be converted to negation normal form.
Next, we normalize the predicate into negation normal form, which pushes negation
into the leaves, using the rewrites $\lnot (P \land Q) \to \lnot P \lor \lnot Q$,
$\lnot (P \lor Q) \to \lnot P \land \lnot Q$, and $\lnot \lnot P \to P$,
and $\lnot (P = Q) \to P \neq Q$ (and similar negations for the other atomic predicates),
which we are able to do since our atomic predicates are closed under negation.
Next, $P = P_1 \land P_2$, one can trivially distribute the quantification
over the conjunction, since $\forall x, P_1(x) \land P_2(x)$ is logically equivalent
to $\forall x, P_1(x) \land \forall y, P_2(y)$. This allows us to solve smaller subproblems.

\paragraph{Specialized Preprocessing for $k$-induction}
The $k$-induction solver has a smaller core language, and it reduces unsigned and signed inequalities
into properties that are checked on the \msb{} of the inputs and outputs.
See that $a <_u b$ if and only if upon computing $a - b$, the high borrow bit
is $1$. Therefore, one can build the stream of borrows from the subtraction circuit,
and define $\denote{a <_u b}[i] \equiv \mathsf{subBorrow}(a, b)[i]$, where $\mathsf{subBorrow}$
builds the borrow circuit. Similarly, note that $a <_s b$ if and only if $(a - b).\toInt < 0$.
At any finite bitwidth $w$, we have that $(a - b)|_w.\toInt < 0$ iff $(a - b)|_w.\msb = \true$.
So we define $\denote{a <_s b} \equiv \lnot (a - b)$, since the bit at index $w$  of $\lnot (a - b)$ is $\false$
iff the \msb of $a - b$ is \true, which implies that $(a - b).\toInt$ is negative.

\paragraph{Width Constraints}
Furthermore, we can impose constraints on the width of the bitvector,
by building bitstreams that start at $1$ and become $0$ after $k$ steps, or vice versa.
This encoding provides constraints of the form $w = k$, $w \geq k$, $w \leq k$,
thereby embedding the finite with fragment into the infinite width theory.
The idea is that the predicate $w = k$ corresponds to a bitstream $\langle0\dots010\dots\rangle$,
which has a $1$ at the index $k$, and $0$ everywhere else. Thus, the predicate
will be true at bitwidth $k$, and $0$ everywhere else.
Similarly, we encode $w \leq k$ as a bitstream that is $1$ at all indices less than $k$,
and $0$ afterward, and $w \geq k$ as a bitstream that is $1$ at all indices greater than $k$.
These constraints are important to rule out vacuous cases which occur at width $0$,
and to prove special cases for small bitwidth.
For example, the statement $\forall (w : \N)~(x~y : \BV{w}), w = 1 \to x + y = x \bitxor y$
is only true at width $1$.

\paragraph{Symbolic Normalization.} We run \bvnormalize{}, a custom bitvector
normalizer which rewrites bitvector expressions into simpler bitvector expressions,
which produced smaller problems that are easier to then reflect into our core language.

\section{A Specialized Solver for Mixed Boolean-Arithmetic}
\label{sec:swbv:mba}

We consider a restricted fragment of the larger bitvector theory, called as mixed-boolean-arithmetic \cite{liu2021mba}
for which an efficient reduction to a decidable integer problem exists.
This works on the fragment of the full theory where we allow linear combinations of \emph{boolean} terms,
where boolean terms are terms comprising of purely bitwise operations on atoms (\cref{fig:swbv:mba-syntax}).
This is an interesting fragment, which subsumes both ring properties and bitwise properties,
and allows for mild mixtures of these properties.

\begin{figure}[htb]
\centering
{\footnotesize
\begin{tabular}{llll}
$\mathsf{A}$ & $::=$ & $\mathsf{Const}$ & bitvector constant \\
    & $\mid$ & $\mathsf{Var}$ & bitvector variable \\[4pt]
$\mathsf{B}$ & $::=$ & $\mathsf{A}$ & atom \\
    & $\mid$ & $\mathsf{B} \boxplus \mathsf{B}$ & bitwise combination, $\boxplus \in \{\bitor, \bitand\}$ \\
    & $\mid$ & $\bitnot \mathsf{B}$ & bitwise complement \\[4pt]
$\mathsf{T}$ & $::=$ & $\mathsf{B}$ & boolean term \\
    & $\mid$ & $\mathsf{T} \otimes \mathsf{T}$ & arithmetic combination, $\otimes \in \{+, -\}$ \\[4pt]
$\mathsf{P}$ & $::=$ & $(\mathsf{T} = 0)$ & MBA equation \\
\end{tabular}
}
\caption{Syntax of the theory decided by the MBA solver: boolean
terms~$\mathsf{B}$, built from atoms~$\mathsf{A}$ by purely bitwise
operations, are combined linearly by arithmetic into terms~$\mathsf{T}$ and
equated to zero.}
\label{fig:swbv:mba-syntax}
\end{figure}

The key idea of such a solver is to first prove MBA identities at the one-bit level,
and to show that if they hold at the one-bit level, then they hold at all bitwidths.
We sketch the idea with an example.
Consider the equation $~\forall~{w : \N}~(x~y : \BV{w}),x + y - 2 (x \bitand y) - (x \bitxor y) =_{\bv{w}} 0$.
First, note that to show this equation, it suffices to show the equation interpreted \emph{over the integers},
with $\forall (w : \N)~(x~y : \BV{w}), \Zcal{x} + \Zcal{y} - 2 (\Zcal{x} \bitand \Zcal{y}) - (\Zcal{x} \bitxor \Zcal{y}) =_{\Z} 0$,
where $\Zcal{\cdot} : \BV{w} \to \Z$ is the interpretation of the bitvector as a 2s complement \emph{unsigned} integer.
Intuitively, this is true because an integer is an infinite width bitvector,
and the equation holds for all bitwidths if it holds for the infinite width bitvector.
To prove the example, we induct the width $w$.
For $w = 1$, we prove the equation by exhaustive enumeration:

\begin{tabular}{ccc}
\rowcolor{gray!10}
  x & y & $\Zcal{x} + \Zcal{y} - 2 \Zcal{x} \bitand \Zcal{y} - \Zcal{x} \bitxor \Zcal{y}$ \\
  0 & 0 & $0 + 0 - 2 \cdot (0) - 0 = 0$ \\
\rowcolor{gray!10}
  0 & 1 & $0 + 1 - 2 \cdot (0) - 1 = 0$ \\
  1 & 0 & $1 + 0 - 2 \cdot (0) - 1 = 0$ \\
\rowcolor{gray!10}
  1 & 1 & $1 + 1 - 2 \cdot (1) - 0 = 0$
\end{tabular}

Which shows that the equation holds for $w = 1$.
Next, assume it holds for $w = n$, and for $w = n + 1$,
use the key property that $\Zcal{\cdot}$ \emph{distributes} over bit concatenation and bit operations:
$$\Zcal{\bvfst{x} : \bvsnd{xs}} \bitand \Zcal{\bvfst{y} : \bvsnd{ys}} = \Zcal{\bvfst{x \bitand y} : \bvsnd{xs \bitand ys}},$$
and similarly for all other bitwise operations $\boxplus$.
Also, we compute $\Zcal{\bvfst{x} : \bvsnd{xs}}$ as $2^n \cdot \Zcal{\bvfst{x}} + \Zcal{\bvsnd{xs}}$.
This allows us to perform the following sequence of rewrites on the goal of the induction, to show
that at width $n + 1$, the equation is zero:

{\footnotesize
\begin{align*}
  &\Zcal{\bvpair{x}{xs}} + \Zcal{\bvpair{y}{ys}} - 2 \Zcal{\bvpair{x}{xs}} \bitand \Zcal{\bvpair{y}{ys}} - \Zcal{\bvpair{x}{xs}} \bitxor \Zcal{\bvpair{y}{ys}} \\
  &= \Zcal{\bvpair{x}{xs}} + \Zcal{\bvpair{y}{ys}} - 2 \Zcal{\bvfst{x \bitand y} : \bvsnd{xs \bitand ys}} - \Zcal {\bvfst{x \bitxor y} : \bvsnd{xs \bitxor ys}}~ (\text{distribute})\\
  &= (2^n \cdot \Zcal{\bvfst{x}} + \Zcal{\bvsnd{xs}}) + (2^n \cdot \Zcal{\bvfst{y}} + \Zcal{\bvsnd{ys}}) - \\
  & \qquad 2 (2^n \cdot \Zcal{\bvfst{x \bitand y}} + \Zcal{\bvsnd{xs \bitand ys}}) - (2^n \cdot \Zcal{\bvfst{x \bitxor y}} + \Zcal{\bvsnd{xs \bitxor ys}})~(\text{evaluate})\\
  &= 2^n \cdot (\Zcal{\bvfst{x}} + \Zcal{\bvfst{y}} - 2 \Zcal{\bvfst{x \bitand y}} - \Zcal{\bvfst{x \bitxor y}}) + \\
  & \quad (\Zcal{\bvsnd{xs}} + \Zcal{\bvsnd{ys}} - 2 \Zcal{\bvsnd{xs \bitand ys}} - \Zcal{\bvsnd{xs \bitxor ys}})~(\text{regroup}) \\
  &= 2^n \cdot 0~(\text{by theorem for width 1}) + 0~(\text{by IH}) = 0,
\end{align*}
}

which proves the identity for all bitwidths.

\subsection{Mechanization in Lean}

In our mechanization, we use the above procedure, appropriately generalized for all boolean
operations $\boxplus$, to check that the identity holds at width $1$ by exhaustive enumeration, as the table does.
We use proof by reflection, and we reflect the equation into a term language:

\begin{lean}
inductive Factor
| var (n : Nat) -- Variable index
| and (x y : Factor) -- Bitwise &
| ...
| not (x : Factor) -- Bitwise !

structure Term where
  c : Int
  f : Factor

def Eqn := List Term
\end{lean}

These define a factor, which is a bitwise expression or a variable, terms, which are factors
scaled by a constant integer, and equations, which are terms which are added up to create
the final expression. We first run a preprocessing pass, which normalizes inputs to be of
the form $\sum_i c_i f_i$, where $c_i$ is a constant integer and $f_i$ is the fragment of \texttt{Factor}.
Example rewrites include translating $- (x + y) \to (-x) + (-y)$, and $-x \to \langle 0 \rangle - x$.
Once we have reflected our expressions, we use the key theorem:
%
\begin{lean}
theorem Eqn.forall_width_reflect_zero_of_width_one_denote_zero
    (e : Eqn) (w : Nat) (env : List (BitVec w))
    (h : (∀ env1 : EnvFin 1 e.numVars, Eqn.denoteFin e env1 = 0)) :
    Eqn.reflect eqn env = BitVec.ofInt w 0
\end{lean}
%
which states that to establish, for all widths, that the equation \leaninline{eqn} is equal to zero,
it suffices to establish this for all a one-bit environments (as stated by hypothesis \leaninline{h}).
The proof of this theorem follows along the lines of the example, with the reasoning generalized
to regroup any linear combination. This establishes the identity for all bitwidths.
Furthermore, the algorithm, while subtle, is algorithmically simple to implement,
taking only ~\MBATotalLines{} lines of code with no external dependencies other than
the Lean language's core library.

\subsection{Algorithmic Extension to (Un)Signed Inequalities}

We contribute a novel extension of the core MBA philosophy of extending $1$-bit identities to all bitwidths,
and show that this can be used to extend the equational theory to the inequational theory.
The key tool will be to exploit the assumption that signed and unsigned overflow do not occur.
Compiler IRs such as LLVM express this assumption by being able to define signed (and unsigned) overflow as \emph{undefined behaviour}.
This is necessary, for example, to express the C language semantics, where integer overflow is undefined behavior,
and in LLVM IR, this is expressed using the \texttt{nsw} and \texttt{nuw} flags (no (un)signed wrap) on arithmetic operations.
If we add the assumption that overflow does not occur, then it is indeed the case that $\denote{x + y}_Z = \denote{x}_Z + \denote{y}_Z$,
which enables us to prove bounds on the integer values of the bitvectors, and to have these bounds transfer to the bitvectors themselves.

For example, consider the expression $(x \bitor y) \geq_u x$.
We can show that this is true by showing that $(x \bitor y) - x \geq 0$ is true for $w = 1$ by exhaustive enumeration,
and by a similar regrouping argument as for equalities, show that it holds for all bitwidths.

\section{Evaluation}
\label{sec:swbv:evaluation}

To show that our MBA decision procedure is complete for its fragment,
and to show that our width-independent decision procedures are able to prove real-world identities from compiler development,
we evaluate our decision procedures quantitatively, by measuring the number of problems solved,
as well as with survival plots~\cite{brain2017benchmarking}, to characterize performance trade-offs of the different solvers.

\begin{figure*}[htb]
{\small
\resizebox{\textwidth}{!}{%
\begin{tabular}{rcccc}
  Solver & Hacker's Del (31) & Alive (104) &  MBA-Blast (2500) & \texttt{InstCombine} (\InstCombineTotalProblems{}) \\
\textbf{Width Indep.} & & & & \\
\rowcolor{gray!20}
  \texttt{k-ind (cadical)} & 31 & 54 & \MBAKInductionVerifiedNumSolved{} & \textbf{\InstCombineNormCircuitVerifiedNumSolved{}} \\
\rowcolor{gray!10}
\texttt{automata} & 31 & 54 & \MBAPresburgerNumSolved{} & \InstCombineNormPresburgerNumSolved{} \\
\textbf{Fixed Width} & & & & \\
\rowcolor{gray!10}
  \texttt{bv\_decide} ($w = 64$) & 31 & 54 & \MBABvDecideNumSolved{} & \InstCombineBvDecideNumSolved{}
\end{tabular}
}
}
\caption{Both k-induction and automata based solvers are competitive, and solve similar fragments.
  We provide a fixed-width comparison (which \textbf{does not solve} the infinite width fragment)
  as a coarse measure of the complexity of the problems.}
\label{fig:swbv:table-eval}
\end{figure*}

\noindent \paragraph{Hacker's Delight~(31 problems).}
For a stock of bit-twiddling algorithms, we take the first two chapters of Hacker's Delight \cite{warren2013hacker},
we take the problems that have been translated into Lean (by \texttt{lean-mlir}~\cite{leanmlir}).

\noindent \paragraph{Alive~(104 problems).} We take the test suite of Alive \cite{lopes2015provably}, which is a translation validator for LLVM IR.
We use problems on addition, subtraction, and bitwise operations which have been translated into Lean (by \texttt{lean-mlir}~\cite{leanmlir}),
and evaluate our algorithms on this dataset.
\noindent \paragraph{InstCombine~(\InstCombineTotalProblems{} problems).}
We use \texttt{InstCombine} a dataset of \InstCombineTotalProblems{} bitvector
lemmas, which are
generated from a dataset of 4556 peephole rewrites from LLVM's peephole optimizer, by
reducing the validity of rewrites on LLVM values to equalities of bitvectors.
%
Note that these problems are for fixed bitwidths.
%
Thus, we can run \bvdecide{} on these problems. To run our decision procedures on this dataset,
we treat the fixed width problem as an infinite width problem, and run our decision procedures on them.
This can introduce spurious counterexamples, such as the optimization $x + y = x \oplus y$, which is untrue for bitwidths greater than 1.
Regardless, we are interested in the performance of our decision procedures to prove real-world rewrites,
and thus this is a valuable dataset.

\subsection{Choosing the Right Solver for the Right Problem}
\label{sec:swbv:right-solver}

In order to understand the trade-offs of the different solvers,
we plot the number of problems solved against cumulative wall-clock time,
with problems sorted in ascending order of wall-clock time.
Such a survival plot~\cite{brain2017benchmarking} allows one to quickly read off the total number of problem solved,
as well as the behavior of the different solvers given a fixed amount of time.

\begin{figure*}
  \begin{subfigure}{\textwidth}
  \includegraphics[width=\textwidth]{chapters/single-width-bv/plots/dataset2-cactus-plot.pdf}
  \caption{Survival plot on the MBA-blast dataset.
  See that the specialized MBA solver is an order of magnitude faster than both
    k-induction and the NFA-based solver,
  this accurately captures the trade-offs available in this space, where specialized solvers are fastest,
  and that consuming certificates from fast, external, unverified solvers effectively does scales to moderate problem sizes.}
  \label{fig:swbv:mba-cactus}
  \end{subfigure}
  \begin{subfigure}{\textwidth}
  \includegraphics[width=\textwidth]{chapters/single-width-bv/plots/automata-automata-circuit-cactus-plot.pdf}
  \caption{Survival plot on the InstCombine (LLVM Peephole Rewrites) dataset.
    We treat fixed-width problems as infinite width problems to be solved by our k-induction and automata solvers.
    Since the problem sizes are small, the lean-based automata solver is faster than the k-induction solver
    which repeatedly invokes CaDiCaL for a proof.}
  \label{fig:swbv:instcombine-cactus}
  \end{subfigure}
  \caption{Survival plots for the MBA-blast and InstCombine datasets.}
\end{figure*}

\noindent \paragraph{MBA~(\MBATotalProblems{} problems).} Out of the \MBATotalProblems{} problems (\cref{fig:swbv:mba-cactus}), the k-induction solver solves \MBAKInductionVerifiedNumSolved{} problems out of \MBATotalProblems{},
while the other solvers solve \MBAPresburgerNumSolved{} problems.
The k-induction solver builds circuits that are quadratic in size,
due to the need to build a circuit that forbids loops, ensuring that all paths are simple paths.
This is not too surprising, as the MBA dataset was designed to be a fragment used for obfuscation,
which is hard for SAT solvers to solve.
The NFA-based solver, which takes \MBAPresburgerTotalTime{} total time,
is in the same order of magnitude as the bitblaster
(where the bitblaster solves for fixed $w = 64$, which we provide as a comparison, but \emph{does not} solve the infinite width problem),
where the bitblaster takes \MBABvDecideTotalTime{} total time.
The MBA solver is the fastest, with a geomean time of \MBAMBAGeoMean{},
which is an order of magnitude faster in comparison to both k-induction (\MBAKInductionVerifiedGeoMean{}) and the NFA-based
solver (\MBAPresburgerGeoMean).

The solvers introduced in this chapter, and especially the automata-based one
which executes entirely inside of Lean, would run faster if native compilation
was enabled (by a factor of at least two for the automata-base one, from
experiments with earlier versions of Mathlib and Lean). However, native
compilation does not work with the recent versions of Mathlib and Lean we tested
with, and so our decision procedures are run in interpreted mode in these
benchmarks.

\paragraph{InstCombine (\InstCombineTotalProblems{} problems).} In the \InstCombineTotalProblems{} problems (\cref{fig:swbv:instcombine-cactus}),
\bvdecide{} is effective, since the problem is always fixed to a particular bitwidth,
and solves \InstCombineBvDecideNumSolved{} problems, a large majority.
$k$-induction solves \InstCombineNormCircuitVerifiedNumSolved{} problems (the largest amongst the symbolic width solvers),
and the automata solver solves \InstCombineNormPresburgerNumSolved{} problems.
Interestingly, the automata solver (with preprocessing) is faster than the k-induction solver (with preprocessing) for the LLVM dataset.\footnote{We conjecture that this is because the problems are small and the overhead of invoking a SAT solver is larger than the time it takes to solve the problem.}
The geomean time for the automata solver is \InstCombineNormPresburgerGeoMean{}, while the geomean time for the k-induction solver is \InstCombineNormCircuitVerifiedGeoMean{}. In contrast, bitblasting is slower, with a geomean time of \InstCombineBvDecideGeoMean{}.
Overall, this shows us that the fragments that our arbitrary-width decision procedures tackle occur in a non-trivial fraction of the real-world problems from the LLVM ecosystem,
and that our algorithmically complete procedures for this fragment is a useful addition to the Lean ecosystem.

\paragraph{Root Causes of LLVM Peephole Failures}
We catalog the failures on the LLVM peephole optimizer dataset,
by logging all errors emitted from the preprocessor for reflection,
and manually inspect the failures, using regular expressions to categorize them.
We find the following categories: (1) Rewrites which ought to hold, but are classified as untrue (447 problems).
These occur since we translate fixed-width problems into infinite width-problems, and such example use constants such as $7$,
which should be generalized to e.g. $2^w - 1$; a promising path would be to use
tools such as Hydra~\citep{hydra} to automatically do the
generalization.
(2) Rewrites that involve quantification over multiple widths $w_1, \dots, w_k$, rather than a fixed width (1925 problems).
These rewrites then perform truncation / sign extension of the bitvectors involved.
(3) Theorems whose goal state lie outside our fragment (463 problems). This includes two kinds of theorems:
(3a) those with operations that are truly not expressible in our fragment, such as signed division,
and (3b) those with operations that need further reduction to be expressible in the fragment, such as division-by-constant.
(4) Theorems that involve boolean equalities that lie outside our fragment (565).
For example, a statement such as \texttt{x : BitVec w ⊢ (255\#w <ᵤ x || 255\#w <ₛ x) = (255\#w <ᵤ x)} uses the
boolean fragment with boolean or, instead of being expressed as a proposition.
This requires further preprocessing, using the newly developed \texttt{bool\_to\_prop} tactic in Lean,
which we plan to integrate into our decision procedures.
(5) Unknown boolean disequality (81 problems), which establish disequality between operations such as \texttt{msb}
which are not supported by our preprocessor.
%
As many of these issues are not fundamental problems with our approach, we
believe our decision procedures will soon be able to prove substantially more
test cases.



\section{Related Work: Automata, Presburger, and Parametric Bitvectors}
\label{sec:mono:related}

The introduction (\cref{ch:intro}) placed this thesis among the approaches that bridge
automated speed and kernel-checked trust; here we compare against the technical
antecedents specific to width-independent reasoning.

\paragraph{Bitwidth Independent Boolean Circuits.}
Bitwidth independent circuits were introduced by Pichora~\cite{pichora2003automated}.
Algorithms for deciding this fragment based on unification~\cite{bjorner1998deciding},
families of boolean circuits to establish base cases and inductive cases~\citep{gupta1993parametric,gupta1993representation}
have been explored in the literature.

\paragraph{Automata-Based Presburger Arithmetic Solvers.}
The idea that a Presburger formula can be represented by a finite-state automaton is classical \cite{Büchi1960}.
Moreover, it can be seen that Presburger arithmetic is expressible into WS1S (weak monadic second order logic with one successor),
which can also be decided via automata.
%
\citet{automata-strike-back} use similar techniques in an uncertified context to
decide Presburger arithmetic with bitwise operations: a formula is translated
to an alternating automaton, whose emptiness is checked using $k$-induction and
a SAT solver.
%
While they mostly focus on unbounded natural numbers with bitwise operations
rather than bitvectors (and do not support left shifts, for example), their
high-level approach is similar to ours.

\paragraph{Model Checking Finite Transition Systems.}
Bounded Model Checking allows properties to be proved for upto $n$ states by unrolling the model.
While the method is incomplete, the power of SAT solving meant that BMC gained popularity in the 1990s,
when SAT solvers started outperforming BDDs \cite{biere1999symbolic}.
\emph{$k$-Induction} \cite{sheeran2000checking} establishes that if the past $k$ states are good, then the $k+1$th state is also good,
for all states.
\emph{Interpolation based methods} instead try to prove the safety property
by finding a stronger property $P$ such that it
(i) implies the safety property, and
(ii) if it holds for the previous $k$ states, then it also holds for the next state.
Such properties are discovered by computing Craig interpolants \cite{craig1957linear}.
Such interpolants can be extracted from resolution proofs in a SAT solver \cite{mcmillan2003interpolation}.
\emph{IC3} \cite{bradley2011sat} and related algorithms all produce a sequence of lemmas,
which are inductive relative to previous lemmas and stepwise assumptions.
At convergence, a subset of the generated lemmas imply the safety property to be proven.
Model checkers such as Pono~\cite{mann2021pono}, AVR~\cite{goel2020avr} can use $k$-induction,
IC3, and interpolation to establish inductive invariants.

\paragraph{Mechanized Model Checking Algorithms \& Finite State Machine Libraries.}
The cava LTL model checker~\cite{gregoire2005proving} is formally verified in Isabelle,
and mechanizes executable B\"uchi automata.
This contrasts to our work, where we mechanize \emph{finite} automata,
and consume finite traces, instead of infinite traces.
%
More importantly, our decision procedures, in addition to being verified, can be
used internally to the proof assistant in the form of tactics and generate
proofs without additional assumptions.
%
\citet{kat-pous} developed a tactic in Coq to decide statements in the theory of
Kleene Algebras with Tests (KAT) by translating from to an automaton and
checking for language inclusion.
%
In contrast to this work, the translation from KAT, which is a variant of
regular expressions, to automata is more direct, and, as the statements being
decided were very small, the automata library is not designed for performance:
for example, the transitions are represented as closures which grow with the
number of operations.
%
The ACL2 proof assistant internally builds finite state machines from its
Symbolic Vector Expression (SVEX)~\cite{hunt2017industrial}, which is associated to each signal in a Verilog module.
These finite state machines are typically used for symbolic simulation and (bounded) model checking.
In contrast, we use our finite automata for decision procedures for bitvector arithmetic,
and we are thus interested in the \emph{languages} that our automata consume.
Previous work has abstractly mechanized k-induction~\cite{ishii2020formalizing},
but does not extract executable tactics from their mechanization,
as their mechanization is entirely abstract with no implementation.

\paragraph{Finite Automata Libraries.}
Mata \cite{chocholaty2024mata} is a high-performance automata library implemented in C++, optimized for fast manipulation and analysis of finite automata. It supports efficient automata-theoretic operations,
Mona \cite{henriksen1995mona} is a C-based tool focused on deterministic finite automata, and is tailored for logics like WS1S. Mona, like our approach, compiles logical specifications into minimal DFAs.
LASH \cite{boigelot2004counting} combines automata with arithmetic reasoning, leveraging automata-based representations to solve Presburger arithmetic formulas.
In contrast to our work, none of the above libraries are mechanized, and none are directly integrated into a proof assistant.


\paragraph{Mixed Boolean Arithmetic Solvers.}
The line of work that builds upon the MBA-Blast algorithm \cite{liu2021mba},
which essentially analyzes $1$-bit patterns and extrapolates them to $n$-bits to prove MBA identities.
This was improved upon by SiMBA~\cite{reichenwallner2022efficient} that can deobfuscate all linear MBAs.
These algorithms are specialized, and prove identities in the fragment where we have a linear combination of terms,
where each term is bitwise operations applied to the bitvector atoms.
To our knowledge, our extension of this algorithm to solve the full horn-clause fragment of the theory is novel,
and our work is the first mechanization of these algorithms.
Our modified presentation of the algorithm in terms of linear algebra is also novel, to the best of our knowledge.

\section{Limitations}
\label{sec:mono:limitations}

Three limitations bound the results of this chapter. The first is one of theory:
the procedures decide the \emph{linear-bitwise} fragment of the parametric
theory, and nonlinear operations---multiplication of two unknowns, and the
like---fall outside it. The boundary is not an artefact of our encoding but of
the reduction to automata itself, which relies on each operation being realizable
as a finite-state transduction over the bit positions; multiplication of two
unknowns is not. Deciding the nonlinear fragment requires a genuinely different
tool; we have an unpublished reduction to the first-order theory of the $2$-adic
integers.

The second is one of trust, and we state it plainly because it qualifies the
thesis's own claim. The $k$-induction backend is verified, and every answer it
returns is kernel-checked; but on the hardest unbounded instances we pair it with
rIC3, an \emph{unverified} model checker, because there the unverified tool is
substantially faster. Answers obtained through that path are trusted at the level
of rIC3, not of the Lean kernel. The verified backend remains available for every
such problem---the trade is one of time, not of reach---but it is a trade, and
we record it here as one.

The third concerns what the procedures leave undone rather than what they get
wrong. Of the \InstCombineTotalProblems{} peephole rewrites extracted from LLVM's
test suite, our tools decide about 31\%. Cataloguing the failures
(\cref{sec:swbv:right-solver}) is instructive, because the residue is mostly not
a fragment boundary. The largest single category, by some margin, is rewrites
that quantify over \emph{several} widths at once and truncate or extend between
them---1925 problems, which the mono-width theory cannot even state. A further
447 fail only because our translation from fixed-width problems leaves concrete
constants such as $7$ ungeneralized where $2^w - 1$ was meant, and 565 because
they phrase a goal in the boolean fragment that our preprocessor does not yet
lift into a proposition. Only 463 lie genuinely outside the fragment, through
operations such as signed division. The dominant limitation of this chapter is
therefore not its fragment but its \emph{arity in widths}, and it is exactly the
limitation \cref{ch:multi-width} removes.

The procedure of this chapter decides predicates over a \emph{single} symbolic
width. Many real compiler optimizations, however, mix operands of \emph{several}
symbolic widths at once---truncations, zero- and sign-extensions---which the
mono-width theory cannot express. Lifting the decision procedure to that setting
is the subject of \cref{ch:multi-width}, which shows, perhaps surprisingly, that
the many-width problem reduces with only linear blowup back to the mono-width
solver developed here.
